\scanpage{48}

% B01 | head, reduction and decomposition
Cela nous ramène au cas où $G$ est soit complet, soit linéaire. Dans le \deleted{premier} \inserted{deuxième} cas, \deleted{considérons pour tout entier $n$} $G=G_u\times G_{ss}$ (décomposition en partie semi-simple et unipotente) et cette décomposition est \uline{définie sur $k$}.

% B02 | left of underlined descent assertion
\margin{à vérifier}

% B03 | upper middle, cases and torsion
Donc on peut supposer que $G$ est soit unipotent, soit \deleted{semi-simple} formé d'éléments semi-simples, soit complet. Dans les deux derniers cas, soit $U$ l'ens. des éléments d'ordre fini de $G=G^{\bar k}$; alors le groupe $G^{\deleted{\unsure{\bar k|k}}}/U$ est un groupe sans torsion, de plus il est divisible; pour le voir, on note que $G_0$ (comp. connexe de $0$) est divisible (c'est le th. de Weil pour une variété abélienne, et \deleted{bien sûr} trivial pour le cas où $G\simeq\bar k^{*\,n}$) et \deleted{\unsure{donc|th}} par suite $G\simeq G_0+G/G_0$, \deleted{\unsure{si|et}} (abstraction faite des opérations de $\mathfrak g_k$) d'où résulte que $U$ contient le facteur $G/G_0$ et que $G/U$ s'identifie au \deleted{\illegible{1}} quotient de $G_0$ par $U\cap G_0$, donc est divisible aussi).

% B04 | left centre, note in two directions
% slanted
\margin{lemme : dégager cf plus \unsure{bas|loin} ici le fait de passer à la limite inductive est essentiel}

% B05 | below, consequence
Donc $\deleted{\unsure{H}}\underline{H}^i(\mathfrak g,G/U)=0$, \inserted{pour $i>0$} \deleted{tandis que} d'où

% B06 | lower quarter, comparison isomorphism
\[\underline{H}^i(\mathfrak g,G)\simeq H^i(\mathfrak g,U)\qquad\text{pour }i\geq1\]

% B07 | lower left, oblique beside comparison
% slanted
\margin{dégager un \uline{lemme}, [lorsque $G$ n'a pas de sous-groupe linéaire \unsure{connexe|commun} unipotent non nul]}

% B08 | lower, parentheses and transition
(compte tenu que $G\longrightarrow G/U$ est surjectif.) Or, pour \unsure{tout|$\forall$} $m$, soit $U(m)$ le sous-groupe de $U$ formé des éléments annulés par $m$. On a

% B09 | foot, direct limit
\[H^i(\mathfrak g,U)=\lim_{\underset m\longrightarrow} H^i(\mathfrak g,U(m))\]

% archive sheet number: 47

